\chapter{Dying and End of Life}
\index{Dying and End of Life}

\section*{Definition and Overview}
End of life refers to the period leading up to death, which may last weeks, months, or even years for people with progressive illnesses. Care at the end of life aims to ensure that a person is as comfortable as possible, free from pain and distressing symptoms, and able to live with dignity according to their wishes. This chapter explains what to expect as a person approaches death, the principles of palliative care, and the practical, emotional, and legal matters that families may need to address, including advance care planning.

\section*{Epidemiology}
Most people in many countries die of chronic conditions such as cancer, heart disease, dementia, and respiratory disease, often after a period of declining health. A significant proportion of deaths occur in hospitals, although many people would prefer to die at home, and rates of home death vary widely by region and circumstances. Palliative care services care for a growing number of people as the population ages, and the majority of people with a life-limiting illness benefit from palliative care at some point.

\section*{Aetiology and Pathophysiology}
The process of dying varies with the underlying illness, but several physical changes are common in the final days.
\begin{itemize}
    \item \textbf{Reduced energy:} the body conserves energy, and the person sleeps more and spends longer periods drowsy
    \item \textbf{Reduced appetite and thirst:} the body needs less food and fluid, and forcing these can cause discomfort
    \item \textbf{Changes in breathing:} breathing may become irregular, with periods of no breathing; a rattling sound from secretions may occur
    \item \textbf{Changes in circulation:} the hands and feet may become cool and mottled as blood is directed to vital organs
    \item \textbf{Changes in consciousness:} the person may become less responsive, though hearing is often the last sense to be lost
    \item \textbf{Confusion or restlessness:} may occur as the body's systems change
\end{itemize}

\section*{Clinical Features}
The signs that death is approaching are gradual, and families can be reassured that these changes are part of the natural process.
\begin{itemize}
    \item Increasing drowsiness and time spent asleep
    \item Withdrawal and reduced communication
    \item Loss of interest in food and drink
    \item Difficulty swallowing, and refusal of medication
    \item Irregular breathing and periods when breathing stops
    \item A rattling or gurgling sound from the throat as secretions collect
    \item Cool hands, feet, and skin with a mottled appearance
    \item Confusion, restlessness, or agitation in some people
    \item A change in the pattern of pain, requiring adjustment of medicines
\end{itemize}

\section*{Differential Diagnosis}
The changes of dying are usually clear in the context of a life-limiting illness, but other conditions can mimic aspects of the process.
\begin{itemize}
    \item Reversible causes of drowsiness or confusion, such as infection, dehydration, or medication side effects
    \item Pain or discomfort that is undertreated
    \item Depression and anxiety in the person or family
    \item The effects of opioid and sedative medications
\end{itemize}

\section*{When to Seek Medical Care}
The palliative care team should be contacted for uncontrolled pain, breathlessness, agitation, or new distressing symptoms, and for advice at any hour. Families should not hesitate to call the GP, palliative care service, or hospice for support and reassurance.

\textbf{Contact your local emergency services for:}
\begin{itemize}
    \item Sudden severe breathlessness or chest pain
    \item Sudden collapse or loss of consciousness from which the person cannot be roused
    \item Sudden severe bleeding
    \item Seizures or fits
    \item Any sudden change that is distressing and needs immediate assistance
\end{itemize}

\section*{Diagnosis and Investigations}
At the end of life, the emphasis shifts from investigation to comfort and care.
\begin{itemize}
    \item \textbf{Assessment:} regular review of symptoms, comfort, and the person's wishes by the care team
    \item \textbf{Advance care planning:} documenting the person's values, preferences, and treatment limits, including an advance care directive and appointment of a substitute decision-maker
    \item \textbf{Review of medications:} stopping medicines that no longer help and continuing those that relieve symptoms
    \item \textbf{Minimal investigation:} tests are avoided unless they will change care or comfort
\end{itemize}

\section*{Management}
End-of-life care is provided by a multidisciplinary team and focuses on comfort and dignity.
\begin{itemize}
    \item \textbf{Pain management:} regular and as-needed medicines, often including opioids, given by mouth, injection, or other routes
    \item \textbf{Symptom control:} medicines for breathlessness, nausea, agitation, and secretions
    \item \textbf{Positioning and nursing care:} turning, mouth care, and skin care to maintain comfort
    \item \textbf{Emotional and spiritual support:} counselling, chaplaincy, and time for the person to say what matters to them
    \item \textbf{Family support:} practical help, information, and bereavement support before and after death
    \item \textbf{Place of care:} care can be provided at home, in a hospice, or in hospital according to the person's wishes and needs
\end{itemize}

\section*{Complications}
\begin{itemize}
    \item Uncontrolled pain or other distressing symptoms if care is inadequate
    \item Family distress and complicated grief
    \item Financial and practical strain on families
    \item Unmet wishes about place of death or treatment if planning is not done
    \item Delirium or restlessness, which may need medication
\end{itemize}

\section*{Prognosis and Outcomes}
Death from a life-limiting illness is a natural process, and good palliative care allows most people to die peacefully and free of pain. While the timing of death cannot be predicted precisely, the final phase is usually characterised by the gradual changes described above. Families who are prepared, supported, and involved report a less distressing experience, and their bereavement is eased by knowing that the person received good care and that their wishes were honoured.

\section*{Prevention and Self-Care}
\begin{itemize}
    \item Discuss wishes early and document them in an advance care directive
    \item Appoint a substitute decision-maker and tell them what matters to you
    \item Learn about the palliative care services available in your area
    \item Keep important documents accessible to family: will, power of attorney, and advance care directive
    \item Ask questions of the care team at any time; there is no such thing as a silly question
    \item Encourage family members to look after their own rest and nutrition
    \item Speak to a counsellor or chaplain if you find the process overwhelming
\end{itemize}

\section*{Sources and Further Reading}
The following reputable sources provide further information for healthcare students and patients seeking reliable, current advice about dying and end-of-life care.
\begin{itemize}
    \item Palliative Care Australia. \textit{End-of-life care information and support.} https://palliativecare.org.au/
    \item NHS. \textit{End of life care: what to expect.} https://www.nhs.uk/conditions/end-of-life-care/
    \item Mayo Clinic. \textit{Palliative care: approaches to end-of-life care.} https://www.mayoclinic.org/departments-centers/palliative-care/
    \item MedlinePlus. \textit{End of life issues.} https://medlineplus.gov/endoflifeissues.html
    \item Advance Care Planning Australia. \textit{Planning ahead for end of life.} https://www.advancecareplanning.org.au/
    \item Dying to Talk. \textit{Conversation starter resources for end-of-life planning.} https://www.dyingtotalk.org.au/
\end{itemize}
